\chapter{Benchmarking Sampling Methods}
\label{ch:ownwork}

Short chapter opening: what we built and why a common pipeline in which all samplers target the identical posterior, so that differences in results can be attributed to the samplers themselves. One paragraph.

\section{Design of the Calibration Framework}
\label{sec:design}

The sampler-agnostic posterior definition: Prior, LogPrior, LogLikelihood, LogPosterior classes shared by all samplers. Why this design: any sampler only needs a function that evaluates the (unnormalized) log-posterior pointwise, so the Bayesian model is defined once and reused. Supported prior distributions (uniform, normal, truncated normal) and optional noise calibration. Common output contract across samplers (mcmc\_output.npz, trace.npy, corner plot, ArviZ InferenceData) so that downstream diagnostics are identical. Here I can give a scheme of the class structure / data flow (parameters in, log-posterior out; sampler on top).

%%%%%%%%%%%%%%%%%%%%%%%%%%%%%%%%%%%%%%%%%%%%%%%%%%%%%%%%%%%%%%%%%%%%%%%%
\section{Mathematical Framework}
\label{sec:math}

The formulas promised in Chapter 2: posterior proportional to likelihood times prior; Gaussian log-likelihood with (optionally calibrated) noise sigma; log-domain evaluation and why (numerical underflow). For dynesty: the unit-cube prior transform via the inverse CDF (ppf), and why nested sampling needs likelihood and prior separately instead of a combined log-posterior (avoiding double-counting the prior). Only equations that are referenced later.

%%%%%%%%%%%%%%%%%%%%%%%%%%%%%%%%%%%%%%%%%%%%%%%%%%%%%%%%%%%%%%%%%%%%%%%%
\section{Starting With the Basics: Random-Walk Metropolis Hastings (RWMH)}
\label{sec:rwmh-sampler}

Here, RWMCMC is the only package I wrote so I'll briefly explain how I did it with the code blocks. Also, I will explain the MCMC logic from that algorithm, and how it is implemented in the code. I will also explain the tuning parameters and their effects on the sampling process.

%%%%%%%%%%%%%%%%%%%%%%%%%%%%%%%%%%%%%%%%%%%%%%%%%%%%%%%%%%%%
\section{Sampler Integration}
\label{sec:sampler-integration}

One subsection per sampler: how each one plugs into the shared posterior, its main tuning parameters, and how its output is mapped to the common contract.

\subsection{emcee}
\label{sec:integration-emcee}

Here I will explain each sampler in detail (under different headings), how do they work, what is important about them. I will also explain the tuning parameters and their effects on the sampling process.

\subsection{dynesty}
\label{sec:integration-dynesty}

done

\subsection{flowMC}
\label{sec:integration-flowMC}
planned

\subsection{cdream}
\label{sec:integration-cdream}
planned
\subsection{PyMC}
\label{sec:integration-PyMC}
planned
\subsection{Pigeons.jl}
\label{sec:integration-Pigeons-jl}
planned
%%%%%%%%%%%%%%%%%%%%%%%%%%%%%%%%%%%%%%%%%%%%%%%%%%%%%%%%%%%

\section{Implementation Challenges and Solutions}
\label{sec:implementation-challenges}

The practical difficulties encountered while integrating the samplers, and how each was resolved. Pickling and multiprocessing: closures are not picklable (no importable top-level name), so the prior transform was refactored into a top-level class with \_\_call\_\_, instantiated inside Prior, why class instances work where closures fail. Parallelization of dynesty via fork context and the nprocs parameter. The rpy2 / R-based surrogate constraint: the forward model itself cannot be parallelized (single worker). Environment pinning: dynesty from conda-forge, matplotlib pin due to arviz incompatibility. Each challenge: symptom, cause, solution, lesson learned.

% \begin{equation}
%   f(x) = \frac{1}{\sigma\sqrt{2\pi}} \exp\left(-\frac{(x - \mu)^2}{2\sigma^2}\right)
%   \label{eq:example}
% \end{equation}

% For multiple related equations, use the \texttt{align} environment to align them at the equals sign:

% \begin{align}
%   \mu    &= \frac{1}{N} \sum_{i=1}^{N} x_i \\
%   \sigma &= \sqrt{\frac{1}{N} \sum_{i=1}^{N} (x_i - \mu)^2}
% \end{align}

% Refer to equations by number in the text, e.g.\ Equation~\eqref{eq:example}.

% \subsection{Implementation}
% \label{sec:implementation}

% Describe any relevant implementation details — tools, frameworks, datasets, or parameters used.\footnote{Be precise: include version numbers, dataset sizes, and hardware specs where relevant so that your work is reproducible.}

%%%%%%%%%%%%%%%%%%%%%%%%%%%%%%%%%%%%%%%%%%%%%%%%%%%%%%%%%%%%%%%%%%%%%%%%
% \section{Results}
% \label{sec:results}

% Present your results objectively. Use tables and figures to summarise findings, and refer to them explicitly in the text.

% \begin{table}[ht]
%   \centering
%   \begin{tabular}{lcc}
%     \toprule
%     \textbf{Condition} & \textbf{Metric A} & \textbf{Metric B} \\
%     \midrule
%     Baseline  & 0.00 & 0.00 \\
%     Your method & 0.00 & 0.00 \\
%     \bottomrule
%   \end{tabular}
%   \caption{Summary of results. Cells contain mean values.}
%   \label{tab:results}
% \end{table}

%%%%%%%%%%%%%%%%%%%%%%%%%%%%%%%%%%%%%%%%%%%%%%%%%%%%%%%%%%%%%%%%%%%%%%%%
% \section{Discussion}
% \label{sec:discussion}

% Interpret your results in the context of your research questions and hypotheses. What do the findings mean? Where does your approach succeed and where does it fall short?\footnote{Be honest about limitations — a well-reasoned discussion of weaknesses is a sign of scientific maturity.}
